% LearnTrack — DBMS Project Report (modular LaTeX)
% Compile from this folder:  pdflatex DBMS_PROJECT_REPORT.tex  (run twice for TOC)
% Chapters live in: latex/chapters/*.tex
\documentclass[12pt,a4paper]{report}

\usepackage[utf8]{inputenc}
\usepackage[T1]{fontenc}
\usepackage{lmodern}
\usepackage[margin=1in]{geometry}
\usepackage{setspace}
\usepackage{graphicx}
\graphicspath{{figures/}}
\usepackage{booktabs}
\usepackage{longtable}
\usepackage{array}
\usepackage[hidelinks]{hyperref}
\usepackage{titlesec}
\usepackage{titling}
\usepackage{enumitem}
\usepackage{listings}
\usepackage{xcolor}

\setstretch{1.15}
\setlength{\parindent}{0pt}
\setlength{\parskip}{6pt}
\setlist[itemize]{leftmargin=*,itemsep=3pt,topsep=3pt}
\setlist[enumerate]{leftmargin=*,itemsep=3pt,topsep=3pt}

\titleformat{\section}{\large\bfseries\color{black}}{\thesection.}{0.5em}{}
\titleformat{\subsection}{\normalsize\bfseries}{\thesubsection}{0.5em}{}

\lstset{
  basicstyle=\ttfamily\footnotesize,
  breaklines=true,
  columns=flexible,
  keepspaces=true,
  frame=single,
  rulecolor=\color{black!40},
  backgroundcolor=\color{black!3}
}

\hypersetup{
  pdftitle={LearnTrack DBMS Project Report},
  pdfauthor={Muhammad Ibrahim}
}

\begin{document}

% ---------- Title page ----------
\begin{titlepage}
  \centering
  \vspace*{0.5cm}
  {\Large \textbf{Bahria University Islamabad}}\\[4pt]
  {\normalsize Department of Computer Science}\\
  {\normalsize Faculty of Computing and Information Technology}\\[1.0cm]

    \includegraphics[width=3.2cm]{Bahria.jpg}\\[1.2cm]

  {\LARGE \textbf{LearnTrack}}\\[6pt]
  {\Large \textbf{Database Management Systems Project Report}}\\[4pt]
  {\large Online Learning Platform with Learning Analytics}\\[1.0cm]

  \rule{0.8\textwidth}{0.6pt}\\[0.5cm]
  {\large \textbf{Course:} Database Management System}\\[4pt]
  {\large \textbf{Instructor:} Ms Mehwish Pervaiz}\\
  \rule{0.8\textwidth}{0.6pt}\\[1.0cm]

  {\large \textbf{Submitted By}}\\[4pt]
  {\large Muhammad Ibrahim}\\
  {\large 01-135241-035}\\
  {\large BS-IT 5A}\\[6pt]
  {\large [Group Member 2 Name]}\\
  {\large [Group Member 2 Roll Number]}\\[4pt]
  {\large [Group Member 3 Name]}\\
  {\large [Group Member 3 Roll Number]}\\[1.0cm]

  {\large \textbf{Submission Date:} May 14, 2026}

  \vfill
\end{titlepage}
\newpage

% ---------- Abstract (plain language) ----------
\begin{abstract}
\textbf{LearnTrack} is an online learning web application. Students register, join courses, watch lessons, update progress, and take quizzes. Instructors create courses, add lessons and assessments, and view simple analytics (for example: who watches the most, who completes lessons, who scores low on quizzes). Administrators can oversee users and high-level statistics depending on the deployed features.

All important data is stored in \textbf{PostgreSQL} through \textbf{Supabase}. The database uses normal relational design: separate tables for users, courses, enrollments, lessons, events, progress, quizzes, attempts, and answers. Small lookup tables store fixed lists (activity types, content types, enrollment statuses). Rules such as ``one enrollment per student per course'' are enforced with \textbf{primary keys}, \textbf{foreign keys}, and \textbf{unique constraints}. Heavy report queries are precomputed in \textbf{five materialized views} and refreshed with a single database function. \textbf{Triggers} and small \textbf{SQL functions} keep timestamps updated and help keep instructor profiles consistent.

The application has three layers: a \textbf{React (Vite)} website, an \textbf{Express.js} REST API under \texttt{/api/v1}, and PostgreSQL accessed with the Supabase JavaScript client. Users authenticate with \textbf{JWT} tokens; Express checks roles (\texttt{student}, \texttt{instructor}, \texttt{admin}) before running handlers. This report explains the design in straightforward language; Row-Level Security in PostgreSQL is mentioned in the DDL but the main story here is API-level access control.

\end{abstract}

\tableofcontents
\newpage

% ---------- Body chapters (one file per chapter) ----------
% Chapter 1 — Introduction (plain-language edition)

\chapter{Introduction}
\label{ch:intro}

\section{Why this project matters}
Schools and training centres increasingly teach online. That creates a simple question: \textit{what are learners actually doing}? Are they opening lessons? How long do they watch? Do they skip videos? How do they score on short quizzes?

If we store this information in spreadsheets or scattered files, reports soon disagree with reality. Names get duplicated, totals do not match, and nobody trusts the numbers.

A \textbf{database} fixes this by keeping \textbf{one official copy} of each fact. SQL makes it easy to ask questions like ``how many students finished Module 2?'' or ``who spent the least time on videos last week?'' The database also \textbf{blocks nonsense data}: for example, you cannot enroll someone in a course that does not exist, because the database enforces \textbf{foreign keys}.

\section{Problem we solved}
We needed a system that does three things at once:

\begin{enumerate}[label=(\alph*),leftmargin=*]
  \item \textbf{Day-to-day use.} Students register, join courses, watch content, and take quizzes. Instructors create courses and see basic reports. All of this must feel fast and reliable.
  \item \textbf{Detailed behaviour data.} We store \textbf{events} (play, pause, skip, etc.) and \textbf{watch time in seconds}, not just ``completed yes/no.'' That supports honest analytics later.
  \item \textbf{Reports without killing the server.} Counting and joining millions of raw rows on every button click is expensive. We therefore store some \textbf{pre-calculated summaries} (materialized views) that dashboards read quickly.
\end{enumerate}

\section{What we set out to achieve}
\begin{enumerate}[leftmargin=*]
  \item Design a \textbf{clear relational model}: separate tables for users, courses, enrollments, lessons, quizzes, attempts, and small ``lookup'' lists (status names, activity types, etc.).
  \item Put \textbf{business rules in the database} where possible: primary keys, foreign keys, checks, and uniqueness so the app cannot accidentally create duplicate enrollments or orphan rows.
  \item Build \textbf{analytics snapshots} (materialized views) keyed by student and course, so instructor dashboards stay responsive.
  \item Keep \textbf{history when we hide things}: soft deletes for users and courses, timestamps on important rows, and full quiz attempt + answer detail for auditing.
  \item Expose everything through a \textbf{REST API} (\texttt{/api/v1/...}) so the React app never talks SQL directly.
\end{enumerate}

\section{What this report covers (and what it does not)}
\textbf{In scope:} How tables relate to each other; why we chose UUIDs for people and integers for courses; how enrollments, activity, progress, and quizzes map to SQL; how materialized views and refresh functions support analytics; how API routes match database writes.

\textbf{Out of scope:} Machine learning grade prediction; native mobile apps; heavy research on learning psychology. Those could build \textit{on top of} this database later.

\section{How to read this report}
\begin{itemize}[leftmargin=*]
  \item Chapter~\ref{ch:overview} gives the big picture: browser, API, Supabase PostgreSQL.
  \item Chapter~\ref{ch:req} lists requirements and which tables satisfy them.
  \item Chapters~\ref{ch:conceptual} and~\ref{ch:logical} explain the entity model and normalization; placeholders show where your ERD and relational diagram go.
  \item Later chapters cover indexes, the full data dictionary, triggers, transactions, analytics queries, workflows, SQL examples, frontend data flow, testing, deployment, and conclusions.
\end{itemize}

% Chapter 2 — System Overview

\chapter{System Overview}
\label{ch:overview}

\section{One-page summary}
Table~\ref{tab:onepage} gives a quick snapshot anyone can read before the technical sections.

\begin{table}[htbp]
\centering
\caption{LearnTrack at a glance}
\label{tab:onepage}
\begin{tabular}{@{}p{3.2cm}p{9.6cm}@{}}
\toprule
\textbf{Topic} & \textbf{Plain explanation} \\
\midrule
Who uses it? & Students learn; instructors create courses and view analytics; admins manage users and platform stats. \\
Where does data live? & All lasting data is in \textbf{PostgreSQL} on \textbf{Supabase}. The React app never touches the database directly. \\
How does the app talk to the DB? & The \textbf{Express} backend uses the official Supabase JavaScript client (\texttt{supabase.from(...)}, \texttt{.rpc(...)}). \\
How do we know who is logged in? & After login, the client sends a \textbf{JWT} (Bearer token). The API checks the token and the user's \textbf{role} before running handlers. \\
How do dashboards stay fast? & Heavy ``join everything'' reports read from \textbf{materialized views} (\texttt{mv\_*}) that we refresh on a schedule or via an admin action. \\
Where is it hosted? & Typical setup: frontend on \textbf{Vercel}, API on \textbf{Render}, database on \textbf{Supabase}. \\
\bottomrule
\end{tabular}
\end{table}

\section{Three layers (simple view)}
Imagine a stack:

\begin{enumerate}[leftmargin=*]
  \item \textbf{Top — React (Vite).} Buttons, forms, charts. Sends HTTPS requests with JSON bodies.
  \item \textbf{Middle — Express API.} Validates input, checks JWT and role, runs the right sequence of database calls.
  \item \textbf{Bottom — PostgreSQL.} Stores rows, enforces constraints, runs triggers and refresh commands for views.
\end{enumerate}

Nothing else is ``magic'': if the UI shows ``15 enrollments,'' that number came from a \texttt{SELECT} or \texttt{COUNT} the API ran moments ago.

\begin{table}[htbp]
\centering
\caption{Layers and technologies}
\label{tab:layers}
\begin{tabular}{@{}llp{6.8cm}@{}}
\toprule
\textbf{Layer} & \textbf{Technology} & \textbf{What it does} \\
\midrule
Presentation & React + Vite + React Router &
Shows different pages for student, instructor, and admin. Uses Axios with a stored token. \\
Application & Express.js &
Defines routes under \texttt{/api/v1}; verifies JWT; validates forms (\texttt{express-validator}). \\
Database access & Supabase JS client &
Builds queries without raw SQL strings in most controllers; calls PostgreSQL functions when needed. \\
Storage & PostgreSQL &
Stores tables, indexes, triggers, and materialized views; guarantees ACID transactions. \\
\bottomrule
\end{tabular}
\end{table}

\section{API layout}
Every REST path starts with \texttt{/api/v1}. Examples students and instructors actually hit:

\begin{itemize}[leftmargin=*]
  \item \textbf{Auth:} register, login, logout (\texttt{/auth/...}).
  \item \textbf{Courses:} list published courses, get one course with lessons and quizzes, instructor ``my courses,'' create/update content (\texttt{/courses/...}).
  \item \textbf{Enrollments:} join a course, see my enrollments (\texttt{/enrollments/...}).
  \item \textbf{Learning signals:} log activity (\texttt{/activity}), update progress (\texttt{/progress/...}).
  \item \textbf{Quizzes:} create quizzes (instructor), submit attempts (student) (\texttt{/quizzes/...}).
  \item \textbf{Analytics:} active students, completion, skipped lessons, underperforming learners, dashboards (\texttt{/analytics/...}).
\end{itemize}

\section{Roles in plain English}
\begin{description}[style=nextline]
  \item[Student]
  Browses published courses, enrolls, watches lessons, sends progress updates, takes quizzes. Cannot edit other people's courses.

  \item[Instructor]
  Owns courses through the chain \texttt{users} $\rightarrow$ \texttt{instructors} $\rightarrow$ \texttt{courses}. Creates lessons and quizzes. Analytics only show data for \textbf{their} course IDs.

  \item[Admin]
  Can manage users and platform-wide views (depending on implemented routes). Sees aggregate analytics across the whole site where the API allows it.
\end{description}

\section{Security model (how this report describes it)}
This report treats security as \textbf{application-level}: JWT verification, role checks in Express, and queries filtered by owner or enrollment.

The SQL script also defines \textbf{Row Level Security} policies for direct Supabase clients, but the main coursework story here is: \textit{the Node API uses the service role and implements rules in code}. That matches how most routes are written in the project controllers.

% Chapter 3 — Requirements Analysis

\chapter{Requirements Analysis}
\label{ch:req}

\section{Who cares about the database?}
Four groups interact with LearnTrack:

\textbf{Students} want one profile, correct enrollment status, visible progress, and quiz scores they can trust. If two screens show different percentages, the database design failed.

\textbf{Instructors} need to publish ordered lessons, attach quizzes, see who enrolled, and skim engagement (watch time, completion, weak quiz averages).

\textbf{Administrators} need user counts, course counts, and sometimes the ability to deactivate accounts or remove bad data---without breaking foreign keys.

\textbf{Developers and DBMS course evaluators} need a single \texttt{ddl.sql} file that rebuilds the same schema anywhere, with constraints documented so behaviour is predictable.

\section{Functional requirements (mapped to tables)}
Table~\ref{tab:funcreq} links plain-language needs to the tables that satisfy them. Use it during vivas: ``Requirement R5 is activity logging; that is \texttt{activity\_log} plus \texttt{activity\_types}.''

\begin{table}[htbp]
\centering
\caption{What the system must do, and where it lives in the schema}
\label{tab:funcreq}
\small
\begin{tabular}{@{}p{8cm}l@{}}
\toprule
\textbf{Requirement (plain wording)} & \textbf{Main tables} \\
\midrule
R1: Every person has one account and a role (student / instructor / admin). & \texttt{users}; instructors also use \texttt{instructors}. \\
R2: Courses belong to an instructor; we can hide a course without deleting history. & \texttt{courses}, \texttt{instructors} \\
R3: Each course has ordered lessons (video, document, quiz marker). & \texttt{content}, \texttt{content\_types} \\
R4: A student joins a course once; status can change (active, completed, dropped). & \texttt{enrollments}, \texttt{enrollment\_statuses} \\
R5: Store fine-grained events: play, pause, skip, plus seconds watched. & \texttt{activity\_log}, \texttt{activity\_types} \\
R6: Store how far each student got on each lesson (0--100\%). & \texttt{content\_progress} \\
R7: Full quiz engine: questions, options, attempts, per-answer rows. & \texttt{quiz}, \texttt{quiz\_questions}, \texttt{question\_options}, \texttt{quiz\_attempts}, \texttt{quiz\_answers} \\
R8: Dashboards show summaries without recomputing giant joins every time. & Materialized views (\texttt{mv\_*}); controller queries \\
\bottomrule
\end{tabular}
\end{table}

\section{Non-functional requirements (what ``good'' means)}
\begin{itemize}[leftmargin=*]
  \item \textbf{Data integrity.} Foreign keys stop dangling enrollments. Unique constraints stop duplicate enrollments for the same student and course.

  \item \textbf{Fair concurrency.} Two enrollment requests at the same second should not create two rows for the same pair. The database \texttt{UNIQUE (user\_id, course\_id)} handles this even if the application forgets to check.

  \item \textbf{Speed for dashboards.} Instructors expect charts to load in seconds. Pre-joined materialized views trade a bit of staleness for fast reads.

  \item \textbf{Maintainability.} Lookup tables (\texttt{activity\_types}, etc.) mean we rename ``skip'' in one place instead of editing thousands of rows.

  \item \textbf{Reset and demo.} The DDL drops and recreates objects in a safe order so coursework and demos can start from a clean slate (\texttt{ddl.sql} + optional \texttt{seed.sql}).

  \item \textbf{Traceability.} Timestamps (\texttt{created\_at}, \texttt{event\_at}, \texttt{attempt\_date}) let us reconstruct timelines when something looks wrong.

  \item \textbf{Room to grow.} Event tables (\texttt{activity\_log}) can accumulate many rows; indexes on user, content, and time support filtered analytics.
\end{itemize}

\section{Requirements traceability tip}
When writing exams or presentations, pick one user story per requirement: e.g.\ ``Student submits quiz'' touches \texttt{quiz\_attempts}, \texttt{quiz\_answers}, and indirectly \texttt{quiz\_questions}, \texttt{question\_options}. Saying the chain out loud proves you understand the schema.

% Chapter 4 — Conceptual Database Design

\chapter{Conceptual Database Design}
\label{ch:conceptual}

\section{What ``conceptual'' means here}
Before we talk about SQL types and indexes, we draw the world in business terms: \textbf{entities} (things), \textbf{relationships} (how things link), and \textbf{cardinality} (one-to-many, many-to-many). The conceptual model should make sense to a non-programmer: ``A course has many lessons; a student joins many courses through enrollment.''

\section{Four storylines}
\textbf{Teaching side.} A person who can teach has a row in \texttt{instructors}. That profile owns \textbf{courses}. Each course contains many \textbf{content} rows (lessons), ordered for display.

\textbf{Learning side.} A student links to a course through \textbf{enrollment}. That row is the permission slip: without it, we should not attribute activity to ``being in the course'' for analytics.

\textbf{Behaviour side.} While learning, the app writes \textbf{activity\_log} rows (many events per student per lesson) and upserts \textbf{content\_progress} (one summary percentage per student per lesson). Logs give detail; progress gives quick ``how far'' reads.

\textbf{Assessment side.} Each course can have \textbf{quizzes}. A quiz has \textbf{questions} and each question has \textbf{options}. When a student submits, we store one \textbf{attempt} and many \textbf{answer} rows---one per question.

\section{Cardinality cheat sheet}
\begin{center}
\begin{tabular}{@{}lll@{}}
\toprule
\textbf{From} & \textbf{To} & \textbf{Meaning in words} \\
\midrule
\texttt{users} & \texttt{instructors} & At most one instructor profile per user (only teachers/admins). \\
\texttt{instructors} & \texttt{courses} & One instructor teaches many courses. \\
\texttt{users} + \texttt{courses} & \texttt{enrollments} & Many-to-many: resolved by an enrollment row (student + course). \\
\texttt{courses} & \texttt{content} & One course has many lessons. \\
\texttt{users}, \texttt{content} & \texttt{activity\_log} & Many events over time. \\
\texttt{users}, \texttt{content} & \texttt{content\_progress} & At most one progress row per pair (unique constraint). \\
\texttt{courses} & \texttt{quiz} & Many quizzes per course. \\
\texttt{quiz} & \texttt{quiz\_questions} & Many questions per quiz. \\
\texttt{quiz\_questions} & \texttt{question\_options} & Many options (MCQ) or two (true/false). \\
\texttt{users}, \texttt{quiz} & \texttt{quiz\_attempts} & Student can attempt multiple times if allowed. \\
\texttt{quiz\_attempts} & \texttt{quiz\_answers} & One row per question answered in that attempt. \\
\bottomrule
\end{tabular}
\end{center}

\section{Why UUID for people and integers for courses}
\textbf{UUID} for \texttt{users.user\_id} matches Supabase Auth: every signup gets a unique ID that never clashes with another project or environment.

\textbf{SERIAL integers} for \texttt{course\_id}, \texttt{content\_id}, etc.\ keep joins compact and human-readable in demo data (\texttt{course\_id = 1}).

\section{Conceptual Entity--Relationship Diagram (ERD)}
\label{sec:erd-fig}

The ERD shows \textbf{entity names} and \textbf{relationship lines}, usually without every column. Standard tools: draw.io, Lucidchart, dbdiagram.io. Export to PDF or PNG and place it as \texttt{figures/erd.pdf} or \texttt{figures/erd.png}.

\textbf{Minimum entities to draw:} \texttt{users}, \texttt{instructors}, \texttt{courses}, \texttt{enrollments}, \texttt{content}, \texttt{activity\_log}, \texttt{content\_progress}, \texttt{quiz}, \texttt{quiz\_questions}, \texttt{question\_options}, \texttt{quiz\_attempts}, \texttt{quiz\_answers}, plus small lookup entities (\texttt{activity\_types}, \texttt{content\_types}, \texttt{enrollment\_statuses}, \texttt{question\_types}).

\textbf{How to insert your diagram:} Uncomment one \texttt{\textbackslash includegraphics} line below and remove the placeholder box.

\begin{figure}[htbp]
  \centering
  % \includegraphics[width=0.94\textwidth,height=0.55\textheight,keepaspectratio]{erd.pdf}
  % \includegraphics[width=0.94\textwidth,height=0.55\textheight,keepaspectratio]{erd.png}
  \fcolorbox{black!50}{black!2}{%
    \parbox{0.92\textwidth}{%
      \centering\vspace{3.2cm}
      {\large\textbf{Placeholder: Conceptual ERD}}\\[1.2em]
      \textit{Add file} \texttt{latex/figures/erd.pdf} \textit{or} \texttt{erd.png}\\[0.6em]
      \small Entities + relationships + cardinalities
      \vspace{3.2cm}
    }%
  }
  \caption{Conceptual ERD for LearnTrack (replace placeholder with your diagram).}
  \label{fig:erd}
\end{figure}

% Chapter 5 — Logical Design and Normalisation

\chapter{Logical Design and Normalisation}
\label{ch:logical}

\section{From drawing to tables}
Conceptual design answers ``what exists?'' Logical design answers ``how do we split information into tables so updates stay safe and queries stay clear?'' That process is \textbf{normalisation}: remove repeated groups, move repeating lists to their own tables, and connect rows with foreign keys.

\section{Why we use small lookup tables}
Instead of storing the word ``skip'' in every row of \texttt{activity\_log}, we store a tiny integer \texttt{type\_id} that points to \texttt{activity\_types}. Benefits:

\begin{itemize}[leftmargin=*]
  \item If we ever rename a label, we update \textbf{one row} in the lookup table.
  \item Typos cannot create fake categories---only listed types exist.
  \item Storage is smaller when millions of events accumulate.
\end{itemize}

\section{Functional dependencies (informal)}
\begin{itemize}[leftmargin=*]
  \item One \texttt{user\_id} determines current profile fields (name, email, role), with \texttt{email} separately unique.
  \item One \texttt{course\_id} determines its instructor and metadata (title, category), when not deleted.
  \item The pair (\texttt{user\_id}, \texttt{course\_id}) should identify \textbf{at most one} enrollment---enforced by a UNIQUE constraint.
\end{itemize}

\section{Why quiz answers live in their own table}
We could imagine cramming all answers into the attempt row as JSON. Instead, \texttt{quiz\_answers} stores one row per question per attempt. That mirrors the relational model used in textbooks: repeating groups become child rows. It also makes future reporting easier (``which option was chosen most often?'').

\section{Referential actions (delete behaviour)}
\textbf{Cascade} on user-owned child rows (sessions, answers) means removing a user removes dependent junk cleanly---important for privacy requests.

\textbf{Restrict} on lookup tables ensures we do not delete ``active'' status while enrollments still reference it.

\section{Relational Data Model Diagram (RDM)}
\label{sec:rdm-fig}

The RDM shows each \textbf{table name}, \textbf{primary key}, \textbf{foreign keys}, and major columns. Materialized views (\texttt{mv\_*}) can appear dashed or in a separate ``analytics'' panel because they are \textbf{derived}, not base tables.

\textbf{Insert your file} as \texttt{figures/rdm.pdf} or \texttt{figures/rdm.png}; uncomment \texttt{\textbackslash includegraphics} below.

\begin{figure}[htbp]
  \centering
  % \includegraphics[width=0.98\textwidth,height=0.62\textheight,keepaspectratio]{rdm.pdf}
  % \includegraphics[width=0.98\textwidth,height=0.62\textheight,keepaspectratio]{rdm.png}
  \fcolorbox{black!50}{black!2}{%
    \parbox{0.94\textwidth}{%
      \centering\vspace{3.6cm}
      {\large\textbf{Placeholder: Relational schema (RDM)}}\\[1.2em]
      \textit{Add file} \texttt{latex/figures/rdm.pdf} \textit{or} \texttt{rdm.png}\\[0.6em]
      \small Tables with PK/FK arrows; include lookups and quiz subgraph
      \vspace{3.6cm}
    }%
  }
  \caption{Relational data model for LearnTrack (replace placeholder with your diagram).}
  \label{fig:rdm}
\end{figure}

\section{Normalisation takeaway for presentations}
Say this in one breath: ``We separated identities, courses, enrollments, events, progress, and quiz structure into different tables so each fact lives in one place, connected by keys.'' That satisfies most DBMS oral expectations without jargon overload.

% Chapter 6 — Physical Design

\chapter{Physical Design}
\label{ch:physical}

\section{What ``physical'' means}
Physical design is about \textbf{performance on real hardware}: which columns we index, which data types we pick, and how we refresh heavy summaries. Students often confuse this with ``logical'' design---remember: logical is ``what tables exist''; physical is ``how PostgreSQL finds rows fast.''

\section{Indexes we actually rely on}
Indexes are like book indexes: they speed up lookups but cost space and slightly slow down inserts. LearnTrack indexes match how the API filters data:

\begin{itemize}[leftmargin=*]
  \item \textbf{Courses by instructor.} Instructors list ``my courses'' with \texttt{courses.instructor\_id}. Index on \texttt{instructor\_id} avoids scanning every course.

  \item \textbf{Enrollments by course.} Roster and enrollment counts filter \texttt{enrollments.course\_id}.

  \item \textbf{Activity by user and time.} Analytics often asks ``what did this student do recently?'' Indexes on \texttt{activity\_log(user\_id)} and \texttt{activity\_log(event\_at DESC)} support that pattern.

  \item \textbf{Quiz attempts.} Reporting joins attempts by \texttt{user\_id} and \texttt{quiz\_id}; indexes there keep quiz history queries usable as data grows.

  \item \textbf{Partial index on published courses.} Listing the public catalogue only needs rows where \texttt{deleted\_at IS NULL} and often \texttt{is\_published}. A partial index shrinks the scanned set.
\end{itemize}

\section{Materialized views and their indexes}
Materialized views store \textbf{precomputed query results} on disk. We put unique indexes on natural keys such as (\texttt{user\_id}, \texttt{course\_id}) for performance summaries so:

\begin{itemize}[leftmargin=*]
  \item lookups by student+course are fast, and
  \item PostgreSQL can refresh some views \texttt{CONCURRENTLY} in setups that support it (less blocking during refresh).
\end{itemize}

\section{Data types (why they matter)}
\begin{itemize}[leftmargin=*]
  \item \textbf{TIMESTAMPTZ} stores ``when'' in absolute time; clients can display local zones without changing stored facts.

  \item \textbf{NUMERIC(5,2)} for percentages avoids floating-point rounding surprises in grade averages.

  \item \textbf{UUID} for \texttt{user\_id} aligns with Supabase authentication IDs.

  \item \textbf{INT / SERIAL} for course and content IDs keeps joins narrow and simple in demo queries.
\end{itemize}

\section{Trade-off summary}
We accept more disk space and occasional refresh time for materialized views because instructor dashboards would otherwise run expensive multi-join aggregates on every page load. Raw event tables stay authoritative; summaries can lag by minutes unless refreshed---a conscious engineering trade teachers should understand when demoing.

% Chapter 7 — Relational Schema and Data Dictionary

\chapter{Relational Schema and Data Dictionary}
\label{ch:dict}

\section{How to use this chapter}
Each subsection describes one table group. In a viva, you can point to the table name and say \textit{what row represents in real life} and \textit{which other tables it points to}.

\section{Lookup (enumeration) tables}
These tables exist only to give stable codes for short lists:

\begin{description}[style=nextline]
  \item[\texttt{activity\_types}]
  Examples: \texttt{play}, \texttt{pause}, \texttt{skip}, \texttt{seek}, \texttt{complete}. Each row gets \texttt{type\_id}; \texttt{activity\_log} stores that ID instead of repeating text.

  \item[\texttt{content\_types}]
  \texttt{video}, \texttt{document}, \texttt{quiz}---tells the UI what kind of lesson row this is.

  \item[\texttt{enrollment\_statuses}]
  \texttt{active}, \texttt{completed}, \texttt{dropped}. Enrollment rows reference one status at a time.

  \item[\texttt{question\_types}]
  \texttt{mcq}, \texttt{true\_false}. Questions reference a type so grading UI knows what shape to expect.
\end{description}

\section{People and ownership}
\textbf{\texttt{users}} holds everyone: students, instructors, admins. Primary key is UUID \texttt{user\_id}. Important columns: \texttt{full\_name}, unique \texttt{email}, \texttt{role}, optional profile fields, \texttt{is\_active}, soft-delete \texttt{deleted\_at}, timestamps.

\textbf{\texttt{instructors}} extends users who teach. Surrogate key \texttt{instructor\_id} is what \texttt{courses} references---cleaner than repeating long UUIDs on every course row. One-to-one with a subset of users (\texttt{user\_id} unique).

\textbf{\texttt{courses}} stores title, description, category, publish flag, soft delete. Foreign key \texttt{instructor\_id} $\rightarrow$ \texttt{instructors}.

\section{Participation and content}
\textbf{\texttt{enrollments}} is the bridge: which student is in which course, with which status. The pair (\texttt{user\_id}, \texttt{course\_id}) is \textbf{unique}---you cannot enroll twice in the same course as the same person.

\textbf{\texttt{content}} rows are lessons inside a course: ordered by \texttt{sort\_order}, typed via \texttt{content\_types}, may store URL or body text.

\section{Behaviour and progress}
\textbf{\texttt{activity\_log}} is the ``fine detail'' table: who, which lesson, which event type, how many seconds watched (\texttt{watch\_time}), when (\texttt{event\_at}). Analytics like ``top watchers'' aggregate this table.

\textbf{\texttt{content\_progress}} stores one rolling percentage per (\texttt{user\_id}, \texttt{content\_id}). The app upserts here when the student watches so the UI can show progress bars without scanning all historical events.

\section{Quiz subsystem (five tables)}
\textbf{\texttt{quiz}} --- one row per assessment attached to a \texttt{course\_id}: title, time limit, pass score, allow multiple attempts, publish flag.

\textbf{\texttt{quiz\_questions}} --- question text, points, order, type.

\textbf{\texttt{question\_options}} --- answer choices and \texttt{is\_correct} for automatic grading.

\textbf{\texttt{quiz\_attempts}} --- one row each time a student submits; stores final \texttt{score}, \texttt{passed}, \texttt{attempt\_date}.

\textbf{\texttt{quiz\_answers}} --- one row per question in that attempt: which option was chosen and whether it was correct. Unique on (\texttt{attempt\_id}, \texttt{question\_id}) prevents duplicate answers for the same question.

\section{Sessions}
\textbf{\texttt{user\_sessions}} stores a hash of the JWT (not the raw token), login time, logout time, expiry. The API can revoke sessions by marking logout, supporting ``log out everywhere'' style behaviour.

\section{Materialized views (analytics snapshots)}
These are \textbf{not} hand-edited tables; they are refreshed from base tables.

\textbf{\texttt{mv\_performance\_summary}} --- per student per course: total watch seconds, average quiz score, completion-style percentage.

\textbf{\texttt{mv\_active\_students}} --- platform-wide engagement rollup from activity (used in admin-style paths).

\textbf{\texttt{mv\_underperforming\_students}} --- quiz averages per learner for global lists.

\textbf{\texttt{mv\_skipped\_content}} --- counts skip events per lesson.

\textbf{\texttt{mv\_completion\_rates}} --- average progress percent per learner per course from \texttt{content\_progress}.

\section{Table count}
The DDL defines \textbf{17 base tables} (including lookups and sessions) plus \textbf{5 materialized views}. This inventory matches \texttt{backend/src/db/ddl.sql}.

% Chapter 8 — Database Programming Artefacts

\chapter{Database Programming Artefacts}
\label{ch:objs}

\section{Why not only tables?}
Tables alone cannot enforce every rule. PostgreSQL lets us attach \textbf{functions} and \textbf{triggers} so that certain updates always run: validate instructor rows, stamp \texttt{updated\_at}, refresh analytics in one call.

\section{Triggers (automatic rules)}
\begin{itemize}[leftmargin=*]
  \item \textbf{Instructor role check.} Before inserting into \texttt{instructors}, the database checks that the linked \texttt{users} row really has role instructor or admin and is not soft-deleted. This stops accidental profiles for students.

  \item \textbf{Auto-create instructor profile.} When a user's role becomes instructor/admin, a trigger can insert the matching \texttt{instructors} row so course creation never fails for missing profile IDs.

  \item \textbf{\texttt{updated\_at} stamps.} A generic trigger sets \texttt{updated\_at = NOW()} on listed tables whenever a row updates---consistent audit trail without repeating code in every API route.
\end{itemize}

\section{Functions}
\textbf{\texttt{refresh\_analytics\_views()}} runs \texttt{REFRESH MATERIALIZED VIEW} for all five \texttt{mv\_*} objects in sequence. The admin API exposes an endpoint that calls this via \texttt{rpc}. Without scheduled refresh, dashboards show slightly old numbers---acceptable for coursework if explained.

\textbf{\texttt{refresh\_performance\_summary()}} is a thin wrapper for backward compatibility; it simply calls the bulk refresh function.

\textbf{\texttt{handle\_new\_auth\_user()}} (when wired in Supabase) inserts into \texttt{public.users} when Auth creates a login---keeps application profile and Auth UUID aligned.

\section{Illustrative SQL: performance summary join pattern}
Listing~\ref{lst:perf} shows the \textbf{LEFT JOIN} pattern behind \texttt{mv\_performance\_summary}. Plain-language meaning: start from every enrollment (student + course). Attach optional facts---watch logs, progress, quiz attempts---without dropping enrollments that have no events yet.

\begin{lstlisting}[caption={Core join pattern for per-student per-course metrics (simplified)},label=lst:perf]
SELECT e.user_id, e.course_id,
       COALESCE(SUM(al.watch_time),0)     AS total_watch_sec,
       COALESCE(AVG(qa.score)::NUMERIC,0)   AS avg_quiz_score,
       COALESCE(AVG(cp.progress_percent)::NUMERIC,0) AS completion_pct
FROM enrollments e
LEFT JOIN content c ON c.course_id = e.course_id
LEFT JOIN activity_log al
       ON al.content_id = c.content_id AND al.user_id = e.user_id
LEFT JOIN content_progress cp
       ON cp.content_id = c.content_id AND cp.user_id = e.user_id
LEFT JOIN quiz q ON q.course_id = e.course_id
LEFT JOIN quiz_attempts qa
       ON qa.quiz_id = q.quiz_id AND qa.user_id = e.user_id
GROUP BY e.user_id, e.course_id;
\end{lstlisting}

\textbf{Why LEFT JOIN?} If a student never opened a quiz, inner joins would erase them from the result. LEFT JOIN keeps the enrollment row and turns missing facts into NULL, which \texttt{COALESCE} turns into zero where appropriate.

\section{Talking point}
Say: ``Triggers encode rules once for every client; functions bundle maintenance tasks like refreshing all analytics views in the correct order.''

% Chapter 9 — Transactional Workflows and API Mapping

\chapter{Transactional Workflows \& API Mapping}
\label{ch:tx}

\section{What we mean by ``transaction'' here}
In everyday speech, a transaction is ``a complete action.'' In the database sense, it can be a \texttt{BEGIN ... COMMIT} block; in this project many operations use Supabase single-row inserts that still behave atomically per statement. The important idea for reports is \textbf{order of checks}: validate permissions first, then write facts that constraints protect.

\section{Enrollment workflow (student joins a course)}
\begin{enumerate}[leftmargin=*]
  \item Load the course row. Confirm it is published (for students) and not soft-deleted.

  \item Check whether an enrollment already exists for this (\texttt{user\_id}, \texttt{course\_id}). If yes, return a friendly error instead of hitting the UNIQUE constraint.

  \item Resolve the ``active'' status ID from \texttt{enrollment\_statuses}.

  \item Insert into \texttt{enrollments}.
\end{enumerate}

\textbf{Why uniqueness matters:} Two tabs open at once might both try to enroll. The application may race; the database UNIQUE constraint still prevents duplicate rows---correctness does not depend only on JavaScript logic.

\section{Quiz submission workflow}
\begin{enumerate}[leftmargin=*]
  \item Confirm the student may take this quiz: enrolled (where required), quiz published, optional single-attempt rule enforced by checking existing \texttt{quiz\_attempts}.

  \item Load all questions and options for scoring.

  \item Grade server-side: compare chosen \texttt{option\_id} with \texttt{is\_correct}; weight by question points; compute percentage score; compare with \texttt{pass\_score}.

  \item Insert one \texttt{quiz\_attempts} row, then insert many \texttt{quiz\_answers} rows.

\end{enumerate}

\textbf{Why grade on server:} The browser cannot be trusted to calculate scores---users could edit DevTools. The database stores the authoritative attempt.

\section{Mapping endpoints to tables (examples)}
\begin{itemize}[leftmargin=*]
  \item \texttt{POST /courses/:id/content} $\rightarrow$ insert \texttt{content} (and resolve \texttt{content\_types} by name).

  \item \texttt{POST /activity} $\rightarrow$ resolve \texttt{activity\_types} by name, insert \texttt{activity\_log}.

  \item \texttt{PUT /progress/:contentId} $\rightarrow$ upsert \texttt{content\_progress}.

  \item \texttt{GET /analytics/...} $\rightarrow$ read \texttt{mv\_*} views and/or aggregate \texttt{activity\_log} / \texttt{quiz\_attempts} depending on role and filters (instructor paths often recompute with enrollment filters).
\end{itemize}

\section{Concurrency in one sentence}
``We rely on constraints like UNIQUE for enrollments so two simultaneous requests cannot corrupt state even if the server code is imperfect.''

% Chapter 10 — Queries, Joins, and Analytics Design

\chapter{Queries, Joins \& Analytics Design}
\label{ch:analytics}

\section{Inner vs outer joins (teaching explanation)}
\begin{itemize}[leftmargin=*]
  \item \textbf{INNER JOIN} keeps only rows that match on \textbf{both} sides. Example: list enrolled students with their names---drop users who somehow lack a profile row (should not happen if FKs hold).

  \item \textbf{LEFT JOIN} keeps every row from the left side and attaches matching right-side rows when they exist. Missing matches show NULL. Use this when the left side is ``official truth'' (enrollments) and the right side is optional facts (quiz attempts).

\end{itemize}

\section{Why summaries sometimes ``multiply'' rows}
If you join enrollments to content to activity \textbf{before} grouping incorrectly, one enrollment can fan out into thousands of activity rows. Aggregates like \texttt{SUM(watch\_time)} must happen at the right grain---usually ``sum watch seconds grouped by user,'' not ``count joined rows.'' Materialized views and careful GROUP BY avoid accidental inflation.

\section{Materialized views vs live queries}
\textbf{Materialized views} trade freshness for speed: good for dashboards refreshed every few minutes.

\textbf{Live queries} on \texttt{activity\_log} respect exact date ranges (``last 7 days'') but cost more CPU. The LearnTrack analytics controller mixes strategies: instructors often use filtered live activity for ``active students'' while other tiles read \texttt{mv\_*} snapshots---know this when explaining why two numbers might not match to the second.

\section{Refresh and staleness}
Calling \texttt{refresh\_analytics\_views()} rebuilds all snapshots. Until then, last refresh time on disk may lag behind new quiz attempts. Operational teams schedule cron jobs or manual refresh after bulk imports.

\section{Answer exam questions like this}
``We store raw facts in normalized tables, aggregate heavy joins into materialized views for dashboards, and use indexes aligned with filter columns in the API.''

% Chapter 11 — Detailed Database Workflows

\chapter{Detailed Database Workflows}
\label{ch:workflows}

\section{Workflow A --- Instructor builds a course}
\begin{enumerate}[leftmargin=*]
  \item Register or log in as instructor; ensure \texttt{users} and \texttt{instructors} rows exist.

  \item Create \texttt{courses} row with correct \texttt{instructor\_id}.

  \item Add \texttt{content} rows in teaching order (\texttt{sort\_order}).

  \item Optional: create \texttt{quiz}, then \texttt{quiz\_questions} and \texttt{question\_options}.

  \item Publish when ready (\texttt{is\_published}) so students see the catalogue entry.
\end{enumerate}

\textbf{Narrative line:} ``Authoring walks down the foreign-key chain: instructor owns course owns lessons owns quiz parts.''

\section{Workflow B --- Student learning session}
\begin{enumerate}[leftmargin=*]
  \item Student enrolls $\rightarrow$ \texttt{enrollments}.

  \item Opens a lesson $\rightarrow$ player sends activity events $\rightarrow$ \texttt{activity\_log}.

  \item Progress updates $\rightarrow$ \texttt{content\_progress} upsert.

  \item Takes quiz $\rightarrow$ \texttt{quiz\_attempts} + \texttt{quiz\_answers}.

\end{enumerate}

\textbf{Narrative line:} ``One session touches both fine-grained logs and coarse progress---that supports analytics depth without killing the progress bar query.''

\section{Workflow C --- Refresh analytics}
\begin{enumerate}[leftmargin=*]
  \item Operational tables accumulate changes throughout the day.

  \item Admin or cron calls stored function to refresh materialized views.

  \item Instructor/admin dashboards read compact rows from \texttt{mv\_*}.

\end{enumerate}

\textbf{Narrative line:} ``This is a classic OLTP + reporting hybrid suitable for coursework scale; a larger company might export to a dedicated warehouse later.''

% Chapter 12 — SQL Patterns and Worked Examples

\chapter{SQL Patterns and Worked Examples}
\label{ch:sqlpatterns}

\section{Purpose}
These queries are templates for demos and vivas. Replace \texttt{\$1} with parameters in application code or Supabase filters.

\section{Roster listing}
Listing~\ref{lst:roster}: everyone enrolled in a course with readable status text.

\begin{lstlisting}[caption={Students enrolled in one course (with status label)},label={lst:roster}]
SELECT
  e.enrollment_id,
  u.user_id,
  u.full_name,
  u.email,
  es.status_name,
  e.enrolled_at
FROM enrollments e
JOIN users u
  ON u.user_id = e.user_id
JOIN enrollment_statuses es
  ON es.status_id = e.status_id
WHERE e.course_id = $1
ORDER BY e.enrolled_at DESC;
\end{lstlisting}

\section{Course-level summary (illustrative)}
Listing~\ref{lst:coursesummary}: joins across enrollments, progress paths, and quizzes. \textbf{Note:} real production queries sometimes split these metrics or use views to avoid double-counting when joins multiply rows---use this example to discuss join grain in class.

\begin{lstlisting}[caption={Illustrative per-course aggregates (discuss join grain carefully)},label={lst:coursesummary}]
SELECT
  c.course_id,
  c.title,
  COUNT(DISTINCT e.user_id) AS enrolled_students,
  COALESCE(ROUND(AVG(cp.progress_percent)::numeric, 2), 0) AS avg_progress_pct,
  COALESCE(ROUND(AVG(qa.score)::numeric, 2), 0) AS avg_quiz_score
FROM courses c
LEFT JOIN enrollments e
  ON e.course_id = c.course_id
LEFT JOIN content ct
  ON ct.course_id = c.course_id
LEFT JOIN content_progress cp
  ON cp.content_id = ct.content_id AND cp.user_id = e.user_id
LEFT JOIN quiz q
  ON q.course_id = c.course_id
LEFT JOIN quiz_attempts qa
  ON qa.quiz_id = q.quiz_id AND qa.user_id = e.user_id
WHERE c.deleted_at IS NULL
GROUP BY c.course_id, c.title;
\end{lstlisting}

\section{Fixing SERIAL sequences after imports}
If you bulk-insert explicit IDs, auto-increment sequences may fall behind the actual maximum ID. Listing~\ref{lst:setval} resets the sequence safely.

\begin{lstlisting}[caption={Align SERIAL sequence with current max ID},label={lst:setval}]
SELECT setval(
  pg_get_serial_sequence('public.courses', 'course_id'),
  COALESCE((SELECT MAX(course_id) FROM public.courses), 1),
  true
);
\end{lstlisting}

\section{Presentation tip}
Walk through one listing line by line: FROM/WHERE choose rows; JOIN connects tables; GROUP BY sets aggregation buckets.

% Chapter 13 — Frontend Interaction (Data Perspective)

\chapter{Frontend Interaction (Data Perspective)}
\label{ch:frontend}

\section{Browsers never run SQL}
The React application only speaks HTTP JSON to Express. Students never receive raw database passwords or service keys. That separation is fundamental: \textbf{the API is the gatekeeper}.

\section{Typical data round-trip}
\begin{enumerate}[leftmargin=*]
  \item User logs in $\rightarrow$ receives JWT stored locally (e.g.\ \texttt{localStorage}).

  \item Axios attaches \texttt{Authorization: Bearer <token>} on later requests.

  \item A page like course detail loads course JSON from \texttt{GET /courses/:id}. That response already joins instructor display names via nested selects on the server side.

  \item Video tracking sends periodic \texttt{POST /activity} and \texttt{PUT /progress/...} calls---these become inserts/upserts in PostgreSQL.

\end{enumerate}

\section{Quiz authoring note}
Instructors may import question banks from JSON in the UI; each question becomes repeated API posts (\texttt{POST /quizzes/:id/questions}). That is many small transactions rather than one giant bulk insert---easier for coursework debugging.

\section{Contract guarantees (what we promise)}
\begin{itemize}[leftmargin=*]
  \item Final quiz scores are computed on the server, not trusted from the browser.

  \item Dashboard numbers come from API JSON responses (often fed by materialized views), not from guesswork in React state.

  \item On HTTP 401, the client clears tokens and redirects to login so half-expired sessions do not silently break writes.

\end{itemize}

\section{Failure handling}
Network timeouts and validation errors surface as toast messages or inline errors. From a DBMS angle, the important lesson is: \textbf{failed requests mean no row was committed}---the database stays consistent even when Wi-Fi drops mid-form.

% Chapter 14 — Implementation Quality and Testing

\chapter{Implementation Quality \& Testing}
\label{ch:qa}

\section{Automated tests}
The backend ships with Jest + Supertest-style tests covering authentication, courses, enrollments, quizzes, and analytics endpoints where implemented. Automated tests catch regressions when schema or controllers change.

The frontend uses Vitest + Testing Library for representative UI flows (example: course list smoke tests). Tests do not replace manual QA but shorten feedback loops during assignment deadlines.

\section{Manual verification matrix}
Use Table~\ref{tab:manualqa} during demos. Replace Pass/Fail with your actual results.

\begin{longtable}{@{}p{4.2cm}p{7.8cm}p{2.2cm}@{}}
\caption{Manual verification checklist}\label{tab:manualqa}\\
\toprule
\textbf{Scenario} & \textbf{What ``good'' looks like} & \textbf{Result} \\
\midrule
Fresh \texttt{ddl.sql} run & All tables, indexes, views, triggers create without errors & Pass / Fail \\
Create course + content & Rows appear in \texttt{courses} and \texttt{content}; FK chain valid & Pass / Fail \\
Duplicate enrollment & Second enroll rejected or safely ignored per API contract & Pass / Fail \\
Quiz submit & Score matches manual calculation from correct options & Pass / Fail \\
Refresh analytics & Materialized views update after RPC / admin action & Pass / Fail \\
Login / logout & Tokens invalidated appropriately; protected routes reject missing JWT & Pass / Fail \\
\bottomrule
\end{longtable}

\section{Operational hygiene}
\begin{itemize}[leftmargin=*]
  \item Keep development and production credentials separate (\texttt{.env} never committed).

  \item Log enough on the API to debug failing inserts without printing secrets.

  \item Validate all mutating payloads (\texttt{express-validator}) so malformed JSON cannot crash controllers.

\end{itemize}

\section{Known coursework angle}
Expect examiners to ask: ``What would you test next?'' Reasonable answers: load testing on \texttt{activity\_log}, backup/restore drill, stricter RLS policies if mobile clients hit Supabase directly.

% Chapter 15 — Deployment Considerations

\chapter{Deployment Considerations}
\label{ch:deploy}

\section{Typical hosting split}
LearnTrack is commonly deployed with:

\begin{itemize}[leftmargin=*]
  \item \textbf{Vercel} for the React build (static assets + SPA routing).

  \item \textbf{Render} (or similar) for the Node Express API.

  \item \textbf{Supabase} for PostgreSQL, Auth, and optional storage.

\end{itemize}

\section{Environment variables that must match}
\begin{itemize}[leftmargin=*]
  \item \texttt{VITE\_API\_URL} on the frontend must point to the deployed API \textbf{including} \texttt{/api/v1} if your client code expects that base.

  \item \texttt{CLIENT\_ORIGIN} (or equivalent CORS setting) on the backend must exactly match the deployed frontend URL---including \texttt{https} and no trailing slash mismatches.

  \item Database URL and service role key live only on the server; never embed them in React.

\end{itemize}

\section{Common production mistakes (and fixes)}
\begin{itemize}[leftmargin=*]
  \item \textbf{404 on all routes:} API base URL missing \texttt{/api/v1} segment.

  \item \textbf{CORS errors in browser console:} backend allowed origin does not match Vercel domain.

  \item \textbf{Stale frontend env vars:} Vite bakes \texttt{VITE\_*} at build time---redeploy after changing them.

  \item \textbf{Cold starts on free tiers:} first request after idle may be slow; retry logic or warm-up ping helps demos.

\end{itemize}

\section{Demo checklist}
\begin{enumerate}[leftmargin=*]
  \item Hit \texttt{/api/v1/health} and confirm JSON status.

  \item Log in from deployed UI.

  \item Create a course and one quiz; submit one attempt as a student.

  \item Trigger analytics refresh if instructor tiles look empty.

  \item Confirm browser console has no CORS errors.

\end{enumerate}

% Chapter 16 — Conclusion

\chapter{Conclusion}
\label{ch:conclusion}

\section{What we demonstrated}
LearnTrack is a complete learning platform backed by a relational database. Students generate real telemetry (activity and progress). Instructors author structured courses and quizzes. Administrators can oversee users and global metrics depending on implemented routes.

From a DBMS course perspective, the project shows:

\begin{itemize}[leftmargin=*]
  \item \textbf{Modeling:} entities, relationships, lookup tables, and clear keys.

  \item \textbf{Integrity:} foreign keys, uniqueness, check constraints, soft deletes where history matters.

  \item \textbf{Analytics:} materialized views that precompute expensive joins for dashboards.

  \item \textbf{Programmability:} triggers and functions for timestamps, instructor rules, and bulk refresh.

  \item \textbf{Integration:} a modern stack (React + Express + Supabase) where SQL remains the source of truth.

\end{itemize}

\section{Limitations (honest)}
\begin{itemize}[leftmargin=*]
  \item Materialized views can lag minutes behind raw logs unless refreshed often.

  \item Very large \texttt{activity\_log} tables eventually need partitioning or archival---outside typical coursework scope.

  \item Row Level Security exists in DDL but the report emphasises API-level auth for clarity.

\end{itemize}

\section{Future extensions}
\begin{itemize}[leftmargin=*]
  \item Scheduled refresh jobs with logged timestamps per view.

  \item Certificate generation stored as immutable rows.

  \item Read replicas or ETL export if analytics traffic grows.

\end{itemize}

\section{Closing statement}
The database is not a filing cabinet---it is the contract that keeps learning data trustworthy. LearnTrack shows how normalized tables, constraints, and targeted summaries work together so teaching teams can see engagement without sacrificing correctness.

\begin{table}[htbp]
\centering
\caption{Summary: what the DBMS project delivered}
\begin{tabular}{@{}p{3.7cm}p{8.9cm}@{}}
\toprule
\textbf{Area} & \textbf{Outcome} \\
\midrule
Model & Normalized schema with lookup tables and explicit relationships. \\
Integrity & Constraints prevent duplicates and orphan rows in normal use. \\
Analytics & Materialized views speed up instructor and admin dashboards. \\
Automation & Triggers/functions keep profiles and timestamps consistent. \\
Deployment & End-to-end operation possible on hosted frontend + API + Supabase. \\
\bottomrule
\end{tabular}
\end{table}


% ---------- Bibliography ----------
\begin{thebibliography}{9}
\bibitem{date}
Date, C.J. \textit{Database Design \& Relational Theory: Normal Forms \& All That Jazz}. (Normal forms and integrity concepts.)
\bibitem{elmasri}
Elmasri \& Navathe. \textit{Fundamentals of Database Systems}. (ER modelling and SQL patterns.)
\bibitem{pg}
PostgreSQL Global Development Group.\ \textit{PostgreSQL Documentation}.\ 
\url{https://www.postgresql.org/docs/}
\bibitem{supa}
Supabase.\ \textit{Supabase Documentation}.\ 
\url{https://supabase.com/docs}
\end{thebibliography}

% Appendices

\appendix

\chapter{REST Resource Inventory (Prefix \texttt{/api/v1})}
\begin{itemize}[leftmargin=*]
  \item \textbf{Health:} \texttt{GET /health} --- simple alive check.

  \item \textbf{Auth:} \texttt{/auth/register}, \texttt{/auth/login}, \texttt{/auth/logout} --- identity and session lifecycle.

  \item \textbf{Users:} profile and admin user management routes as implemented in \texttt{users.js}.

  \item \textbf{Courses:} public catalogue, instructor ``mine,'' single course with nested content and quizzes, CRUD, roster.

  \item \textbf{Enrollments:} list mine, enroll, status updates, admin deletes where enabled.

  \item \textbf{Activity \& progress:} log events, read/update completion percentages.

  \item \textbf{Quizzes:} create/update questions, submit attempts, review attempts.

  \item \textbf{Analytics:} active students, underperforming learners, skipped content, completion rates, instructor course dashboard, admin dashboard, optional materialized view refresh.

\end{itemize}

\chapter{Mini Glossary}
\begin{description}[style=nextline]
  \item[Foreign key]
  A column that must match an existing primary key in another table---keeps relationships honest.

  \item[Materialized view]
  Saved result of a query refreshed on demand or schedule; faster reads, slightly stale data.

  \item[Soft delete]
  Mark a row inactive (\texttt{deleted\_at}) instead of removing it, preserving references.

  \item[Surrogate key]
  Database-generated ID (e.g.\ SERIAL) used purely to identify a row.

  \item[UUID]
  Global unique identifier for users here aligned with Supabase Auth.

\end{description}

\chapter{Repository File Map}
\begin{itemize}[leftmargin=*]
  \item \texttt{backend/src/db/ddl.sql} --- full schema (tables, indexes, views, triggers, RLS).

  \item \texttt{backend/src/db/seed.sql} --- optional demo data.

  \item \texttt{backend/src/db/connection.js} --- Supabase client bootstrap.

  \item \texttt{backend/src/controllers/*.js} --- API logic and queries.

  \item \texttt{learntrack-client/src/api/index.js} --- Axios base URL and endpoint helpers.

  \item \texttt{latex/chapters/*.tex} --- modular chapters for this report (included from \texttt{DBMS\_PROJECT\_REPORT.tex}).

\end{itemize}


\end{document}
